\frfilename{mt121.tex}
\versiondate{21.12.03}
\copyrightdate{1994}

\def\chaptername{Integrasi}
\def\sectionname{Fungsi terukur}

\newsection{121}

Dalam bagian ini, saya mundur sejenak untuk mengembangkan gagasan-gagasan
yang berkaitan dengan aljabar-$\sigma$ himpunan, melanjutkan \S111; di
sini tidak akan disebutkan `ukuran', kecuali dalam latihan. Tujuannya
adalah menetapkan konsep `fungsi terukur' (121C) dan berbagai teknik
yang terkait dengannya.
Contoh tunggal terbaik tentang aljabar-$\sigma$ yang patut diingat ketika
membaca bab ini barangkali adalah aljabar-$\sigma$ subhimpunan Borel dari
$\Bbb R$\cmmnt{ (111G)}; aljabar-$\sigma$ subhimpunan terukur Lebesgue
dari $\Bbb R$\cmmnt{ (114E)} merupakan contoh kedua yang baik.

Sepanjang uraian di sini (mulai dari 121A), saya berusaha menangani
fungsi-fungsi yang tidak didefinisikan pada seluruh ruang $X$ yang
sedang ditinjau. Saya percaya bahwa ada alasan-alasan kuat untuk
menghadapi fungsi semacam itu sejak awal\cmmnt{ (lihat 121G)}; tetapi
tidak dapat disangkal bahwa fungsi-fungsi tersebut menambah kesulitan
teknis. Karena itu, wajar saja bila Anda membaca bab ini sekali dengan
anggapan sementara bahwa semua fungsi terdefinisi di mana-mana, sebelum
kembali untuk menangani kasus umum.


\leader{121A}{Lema} Misalkan $X$ suatu himpunan dan $\Sigma$ suatu
aljabar-$\sigma$ atas subhimpunan-subhimpunan $X$. Misalkan $D$
sembarang subhimpunan $X$ dan tuliskan

\Centerline{$\Sigma_D=\{E\cap D:E\in\Sigma\}$.}

\noindent Maka $\Sigma_D$ adalah aljabar-$\sigma$ atas
subhimpunan-subhimpunan $D$.

\proof{{\bf (i)} $\emptyset=\emptyset\cap D\in\Sigma_D$ karena
$\emptyset\in\Sigma$.

\medskip

{\bf (ii)} Jika $F\in\Sigma_D$, terdapat $E\in\Sigma$ sedemikian
sehingga $F=E\cap D$; sekarang
$D\setminus F=(X\setminus E)\cap D\in\Sigma_D$ karena
$X\setminus E\in\Sigma$.

\medskip

{\bf (iii)} Jika $\langle F_n\rangle_{n\in\Bbb N}$ adalah sembarang
barisan di dalam $\Sigma_D$, maka untuk setiap $n\in\Bbb N$ kita dapat
memilih $E_n\in\Sigma$ sedemikian sehingga $F_n=E_n\cap D$; sekarang
$\bigcup_{n\in\Bbb N}F_n=(\bigcup_{n\in\Bbb N}E_n)\cap D\in\Sigma_D$
karena $\bigcup_{n\in\Bbb N}E_n\in\Sigma$.
}%end of proof of 121A

\medskip

\noindent{\bf Notasi} Saya akan menyebut $\Sigma_D$ sebagai
{\bf aljabar-$\sigma$ subruang} atas subhimpunan-subhimpunan $D$, dan
saya akan mengatakan bahwa anggota-anggotanya {\bf terukur relatif}
di $D$.
$\Sigma_D$ juga kadang-kadang disebut {\bf jejak} $\Sigma$ pada $D$.

\leader{121B}{Proposisi} Misalkan $X$ suatu himpunan, $\Sigma$ suatu
aljabar-$\sigma$ atas subhimpunan-subhimpunan $X$, dan $D$ suatu
subhimpunan $X$. Tuliskan $\Sigma_D$ untuk aljabar-$\sigma$ subruang
atas subhimpunan-subhimpunan $D$. Maka, untuk setiap fungsi
$f:D\to\Bbb R$, pernyataan-pernyataan berikut ekuivalen\cmmnt{, yakni,
jika salah satunya benar maka semuanya benar}:

\quad(i) $\{x:f(x)<a\}\in\Sigma_D$ untuk setiap $a\in\Bbb R$;

\quad(ii) $\{x:f(x)\le a\}\in\Sigma_D$ untuk setiap $a\in\Bbb R$;

\quad(iii) $\{x:f(x)>a\}\in\Sigma_D$ untuk setiap $a\in\Bbb R$;

\quad(iv) $\{x:f(x)\ge a\}\in\Sigma_D$ untuk setiap $a\in\Bbb R$.

\proof{{\bf (i)$\Rightarrow$(ii)} Andaikan (i), dan misalkan
$a\in\Bbb R$. Maka

\Centerline{$\{x:f(x)\le a\}
=\bigcap_{n\in\Bbb N}\{x:f(x)<a+2^{-n}\}\in\Sigma_D$}

\noindent karena $\{x:f(x)<a+2^{-n}\}\in\Sigma_D$ untuk setiap $n$
dan $\Sigma_D$ tertutup terhadap irisan terhitung (111Dd).
Karena $a$ sebarang, (ii) benar.

\wheader{121B}{4}{2}{2}{72pt}

{\bf (ii)$\Rightarrow$(iii)} Andaikan (ii), dan misalkan
$a\in\Bbb R$. Maka

\Centerline{$\{x:f(x)>a\}=D\setminus\{x:f(x)\le a\}\in\Sigma_D$}

\noindent karena $\{x:f(x)\le a\}\in\Sigma_D$ dan $\Sigma_D$
tertutup terhadap pengambilan komplemen. Karena $a$ sebarang, (iii)
benar.

\medskip

{\bf (iii)$\Rightarrow$(iv)} Andaikan (iii), dan misalkan
$a\in\Bbb R$. Maka

\Centerline{$\{x:f(x)\ge a\}
=\bigcap_{n\in\Bbb N}\{x:f(x)>a-2^{-n}\}\in\Sigma_D$}

\noindent karena $\{x:f(x)>a-2^{-n}\}\in\Sigma_D$ untuk setiap $n$
dan $\Sigma_D$ tertutup terhadap irisan terhitung.
Karena $a$ sebarang, (iv) benar.

\medskip

{\bf (iv)$\Rightarrow$(i)} Andaikan (iv), dan misalkan
$a\in\Bbb R$. Maka

\Centerline{$\{x:f(x)<a\}=D\setminus\{x:f(x)\ge a\}\in\Sigma_D$}

\noindent karena $\{x:f(x)\ge a\}\in\Sigma_D$ dan $\Sigma_D$
tertutup terhadap pengambilan komplemen. Karena $a$ sebarang, (i)
benar.
}%end of proof of 121B

\leader{121C}{Definisi} Misalkan $X$ suatu himpunan, $\Sigma$ suatu
aljabar-$\sigma$ atas subhimpunan-subhimpunan $X$, dan $D$ suatu
subhimpunan $X$. Sebuah fungsi $f:D\to\Bbb R$ disebut {\bf terukur}
(atau {\bf $\Sigma$-terukur}) jika memenuhi salah satu, atau secara
ekuivalen semua, syarat (i)-(iv) dalam 121B.

Jika $X$ adalah $\Bbb R$ atau $\BbbR^r$, dan $\Sigma$ adalah
aljabar-$\sigma$ Borel-nya\cmmnt{ (111G)}, suatu fungsi $\Sigma$-terukur
disebut {\bf terukur Borel}. Jika $X$ adalah $\Bbb R$ atau $\BbbR^r$,
dan $\Sigma$ adalah aljabar-$\sigma$ himpunan-himpunan terukur
Lebesgue\cmmnt{ (114E, 115E)}, suatu fungsi $\Sigma$-terukur disebut
{\bf terukur Lebesgue}.

\cmmnt{\medskip

\noindent{\bf Catatan} Tentu saja, kasus utama di sini ialah ketika
$D=X$. Namun, fungsi-fungsi yang terdefinisi sebagian begitu umum, dan
begitu penting, dalam analisis (tinjau, misalnya, fungsi real
$\ln\sin$) sehingga tampaknya layak untuk sejak awal menetapkan
teknik-teknik yang dapat menanganinya secara efisien.

Banyak penulis mengembangkan teori `bilangan real diperluas' pada titik
ini, dengan bekerja pada
$[-\infty,\infty]=\Bbb R\cup\{-\infty,\infty\}$, dan mendefinisikan
keterukuran bagi fungsi-fungsi yang mengambil nilai dalam himpunan ini.
Saya menguraikan garis besar teori semacam itu dalam \S135 di bawah.
}%end of comment

\leader{121D}{Proposisi} Misalkan $X$ adalah $\BbbR^r$ untuk suatu
$r\ge 1$, $D$ suatu subhimpunan $X$, dan $g:D\to\Bbb R$ suatu fungsi.

(a) Jika $g$ terukur Borel, maka fungsi tersebut terukur Lebesgue.

(b) Jika $g$ kontinu, maka fungsi tersebut terukur Borel.

(c) Jika $r=1$ dan $g$ monoton, maka fungsi tersebut terukur Borel.

\proof{{\bf (a)} Hal ini langsung mengikuti definisi dalam 121C,
jika kita mengingat bahwa aljabar-$\sigma$ Borel termuat dalam
aljabar-$\sigma$ Lebesgue (114G, 115G).

\medskip

{\bf (b)} Ambil $a\in\Bbb R$. Tetapkan

\Centerline{$\Cal G=\{G:G\subseteq\BbbR^r\text{ terbuka},
\,g(x)<a\Forall x\in G\cap D\}$,}

\Centerline{$G_0=\bigcup\Cal G
=\{x:\exists\enskip G\in\Cal G,\,x\in G\}$.}

\noindent Maka $G_0$ adalah gabungan himpunan-himpunan terbuka, sehingga
terbuka (1A2Bd). Selanjutnya,

\Centerline{$\{x:g(x)<a\}=G_0\cap D$.}

\noindent\Prf\ {\bf (i)} Jika $g(x)<a$, maka (karena $g$ kontinu)
terdapat $\delta>0$ sedemikian sehingga $|g(y)-g(x)|<a-g(x)$ setiap
kali $y\in D$ dan $\|y-x\|<\delta$. Tetapi
$\{y:\|y-x\|<\delta\}$ terbuka (1A2D), sehingga termasuk dalam
$\Cal G$ dan termuat dalam $G_0$, serta $x\in G_0\cap D$.
{\bf (ii)} Jika $x\in G_0\cap D$, terdapat $G\in\Cal G$ sedemikian
sehingga $x\in G$; sekarang $g(y)<a$ untuk setiap $y\in G\cap D$,
sehingga, khususnya, $g(x)<a$.\ \Qed

Akhirnya, $G_0$, karena terbuka, merupakan himpunan Borel. Karena $a$
sebarang, $g$ terukur Borel.

\medskip

{\bf (c)} Andaikan terlebih dahulu bahwa $g$ tak-menurun. Misalkan
$a\in\Bbb R$ dan tuliskan $E=\{x:g(x)<a\}$. Jika $E=D$ atau
$E=\emptyset$, maka tentu saja himpunan tersebut adalah irisan $D$
dengan suatu himpunan Borel. Jika tidak, $E$ tak kosong dan berbatas atas dalam
$\Bbb R$, sehingga mempunyai supremum $c\in\Bbb R$. Sekarang $E$
haruslah $D\cap\ocint{-\infty,c}$ atau $D\cap\ooint{-\infty,c}$;
bentuk pertama berlaku jika $c\in E$, dan bentuk kedua jika tidak.
Dalam kedua kasus, himpunan tersebut merupakan irisan $D$ dengan suatu
himpunan Borel (lihat 114G).

Demikian pula, jika $g$ tak-menaik, $\{x:g(x)>a\}$ kembali akan menjadi
irisan $D$ dengan salah satu dari $\emptyset$, $\Bbb R$,
$\ocint{-\infty,c}$, atau $\ooint{-\infty,c}$ untuk suatu $c$. Jadi,
dalam kasus ini 121B(iii) akan terpenuhi.

\medskip

\noindent{\bf Catatan} Saya melihat bahwa dalam bagian (b) dari bukti
di atas saya menggunakan beberapa fakta dasar tentang himpunan terbuka
dalam $\BbbR^r$. Fakta-fakta ini dibahas secara terperinci dalam
\S1A2. Jika fakta-fakta tersebut baru bagi Anda, barangkali bijaksana
untuk mengulang argumennya dengan $r=1$, sehingga $D\subseteq\Bbb R$,
sebelum beralih ke kasus umum.
}%end of proof of 121D


\vleader{108pt}{121E}{Teorema}
Misalkan $X$ suatu himpunan dan $\Sigma$ suatu aljabar-$\sigma$
atas subhimpunan-subhimpunan $X$. Misalkan $f$ dan $g$ fungsi bernilai
real yang didefinisikan pada domain $\dom f$, $\dom g\subseteq X$.

(a) Jika $f$ konstan, maka fungsi tersebut terukur.

(b) Jika $f$ dan $g$ terukur, maka $f+g$ juga terukur, dengan
$(f+g)(x)=f(x)+g(x)$ untuk $x\in\dom f\cap\dom g$.

(c) Jika $f$ terukur dan $c\in\Bbb R$, maka $cf$ terukur, dengan
$(cf)(x)=c\cdot f(x)$ untuk $x\in \dom f$.

(d) Jika $f$ dan $g$ terukur, maka $f\times g$ juga terukur, dengan
$(f\times g)(x)=f(x)\times g(x)$ untuk
$x\in \dom f\cap\dom g$.

(e) Jika $f$ dan $g$ terukur, maka $f/g$ juga terukur, dengan
$(f/g)(x)=f(x)/g(x)$ ketika $x\in \dom f\cap\dom g$ dan $g(x)\ne 0$.

(f) Jika $f$ terukur dan $E\subseteq\Bbb R$ suatu himpunan Borel, maka
terdapat $F\in\Sigma$ sedemikian sehingga
$f^{-1}[E]=\{x:f(x)\in E\}$ sama dengan $F\cap\dom f$.

(g) Jika $f$ terukur dan $h$ suatu fungsi terukur Borel dari suatu
subhimpunan $\dom h$ dari $\Bbb R$ ke $\Bbb R$, maka $hf$ terukur,
dengan $(hf)(x)=h(f(x))$ untuk
$x\in\dom(hf)=\{y:y\in\dom f,\,f(y)\in\dom h\}$.

(h) Jika $f$ terukur dan $A$ sembarang himpunan, maka $f\restr A$
terukur, dengan $\dom(f\restr A)=A\cap\dom f$ dan
$(f\restr A)(x)=f(x)$ untuk $x\in A\cap\dom f$.

\proof{ Untuk setiap $D\subseteq X$, tuliskan $\Sigma_D$ untuk
aljabar-$\sigma$ subruang atas subhimpunan-subhimpunan $D$.

\medskip

{\bf (a)} Jika $f(x)=c$ untuk setiap $x\in \dom f$, maka
$\{x:f(x)<a\}=\dom f$ jika $c<a$, dan $\emptyset$ jika tidak; dalam
kedua kasus himpunan tersebut termasuk dalam $\Sigma_{\dom f}$.

\medskip

{\bf (b)} Tuliskan $D=\dom(f+g)=\dom f\cap\dom g$. Jika $a\in\Bbb R$,
tetapkan $K=\{(q,q'):q,\,q'\in\Bbb Q,\,q+q'\le a\}$. Maka $K$ adalah
subhimpunan $\Bbb Q\times\Bbb Q$, sehingga terhitung (111Fb, 1A1E).
Untuk $q\in\Bbb Q$, pilih himpunan-himpunan $F_q$,
$G_q\in\Sigma$ sedemikian sehingga

\Centerline{$\{x:f(x)<q\}=F_q\cap\dom f$,
\quad$\{x:g(x)<q\}=G_q\cap\dom g$.}

\noindent Untuk setiap $(q,q')\in K$, himpunan

\Centerline{$E_{qq'}
=\{x:f(x)<q,\,g(x)<q'\}=F_q\cap G_{q'}\cap D$}

\noindent termasuk dalam $\Sigma_D$.
Akhirnya, jika $(f+g)(x)<a$, kita dapat menemukan
$q\in\ooint{f(x),a-g(x)}$, $q'\in\ocint{g(x),a-q}$, sehingga
$(q,q')\in K$ dan $x\in E_{qq'}$; sedangkan jika $(q,q')\in K$ dan
$x\in E_{qq'}$, maka $(f+g)(x)<q+q'\le a$. Jadi

\Centerline{$\{x:(f+g)(x)<a\}
=\bigcup_{(q,q')\in K}E_{qq'}\in\Sigma_D$}

\noindent menurut 111Fa. Karena $a$ sebarang, $f+g$ terukur.

\medskip

{\bf (c)} Tuliskan $D=\dom f$. Misalkan $a\in\Bbb R$. Jika $c>0$,
maka

\Centerline{$\{x:cf(x)<a\}=\{x:f(x)<\Bover{a}{c}\}\in\Sigma_D$.}

\noindent Jika $c<0$, maka

\Centerline{$\{x:cf(x)<a\}=\{x:f(x)>\Bover{a}{c}\}\in\Sigma_D$.}

\noindent Sedangkan jika $c=0$, maka $\{x:cf(x)<a\}$ adalah $D$ atau
$\emptyset$, seperti dalam (a) di atas, sehingga termasuk dalam
$\Sigma_D$. Karena $a$ sebarang, $cf$ terukur.

\medskip

{\bf (d)} Tuliskan $D=\dom(f\times g)=\dom f\cap\dom g$. Misalkan
$a\in\Bbb R$. Misalkan $K$ adalah

\Centerline{$\{(q_1,q_2,q_3,q_4):q_1,\ldots,q_4\in\Bbb Q,\,uv<a
\text{ setiap kali }u\in\ooint{q_1,q_2},\,v\in\ooint{q_3,q_4}\}$.}

\noindent Maka $K$ terhitung. Untuk $q\in\Bbb Q$, pilih
himpunan-himpunan $F_q$, $F'_q$, $G_q$, $G'_q\in\Sigma$ sedemikian
sehingga

\Centerline{$\{x:f(x)<q\}=F_q\cap\dom f$,
\quad$\{x:f(x)>q\}=F'_q\cap\dom f$,}

\Centerline{$\{x:g(x)<q\}=G_q\cap\dom g$,
\quad$\{x:g(x)>q\}=G'_q\cap\dom g$.}

\noindent Untuk $(q_1,q_2,q_3,q_4)\in K$, tetapkan

$$\eqalign{E_{q_1q_2q_3q_4}
&=\{x:f(x)\in\ooint{q_1,q_2},\,g(x)\in\ooint{q_3,q_4}\}\cr
&=D\cap F'_{q_1}\cap F_{q_2}\cap G'_{q_3}\cap G_{q_4}
\in\Sigma_D;\cr}$$

\noindent maka $E=\bigcup_{(q_1,q_2,q_3,q_4)\in
K}E_{q_1q_2q_3q_4}\in\Sigma_D$.

Sekarang $E=\{x:(f\times g)(x)<a\}$. \Prf\ {\bf (i)} Jika
$(f\times g)(x)<a$, tetapkan $u=f(x)$, $v=g(x)$. Tetapkan

\Centerline{$\eta=\min(1,\Bover{a-uv}{1+|u|+|v|})>0$.}

\noindent Ambil $q_1,\ldots,q_4\in\Bbb Q$ sedemikian sehingga

\Centerline{$u-\eta\le q_1<u<q_2\le u+\eta$,
\quad$v-\eta\le q_3<v<q_4\le v+\eta$.}

\noindent Jika $u'\in\ooint{q_1,q_2}$,
$v'\in\ooint{q_3,q_4}$, maka $|u'-u|<\eta$ dan $|v'-v|<\eta$, sehingga

$$\eqalign{u'v'-uv
&=(u'-u)(v'-v)+(u'-u)v+u(v'-v)\cr
&<\eta^2+\eta|v|+|u|\eta
\le\eta(1+|u|+|v|)
\le a-uv,\cr}$$

\noindent dan $u'v'<a$. Dengan demikian $(q_1,q_2,q_3,q_4)\in K$.
Selain itu, $x\in E_{q_1q_2q_3q_4}$, sehingga $x\in E$. Jadi
$\{x:(f\times g)(x)<a\}\subseteq E$. {\bf (ii)} Di pihak lain, jika
$x\in E$, terdapat $q_1,\ldots,q_4$ sedemikian sehingga
$(q_1,q_2,q_3,q_4)\in K$ dan $x\in E_{q_1q_2q_3q_4}$, sehingga
$f(x)\in\ooint{q_1,q_2}$ dan $g(x)\in\ooint{q_3,q_4}$ serta
$f(x)g(x)<a$. Jadi $E\subseteq\{x:(f\times g)(x)<a\}$.\ \Qed

Dengan demikian $\{x:(f\times g)(x)<a\}\in\Sigma_D$. Karena $a$
sebarang, $f\times g$ terukur.

\medskip

{\bf (e)} Berdasarkan (d), cukup ditunjukkan bahwa $1/g$ terukur.
Sekarang, jika $a>0$,
$\{x:1/g(x)<a\}=\{x:g(x)>1/a\}\cup\{x:g(x)<0\}$; jika $a<0$, maka
$\{x:1/g(x)<a\}=\{x:1/a<g(x)<0\}$; dan jika $a=0$, maka
$\{x:1/g(x)<a\}=\{x:g(x)<0\}$. Semua himpunan ini termasuk dalam
$\Sigma_{\dom 1/g}$.

\medskip

{\bf (f)} Tuliskan $D=\dom f$ dan tinjau himpunan

\Centerline{$\Tau=\{E:E\subseteq\Bbb R,\,f^{-1}[E]\in\Sigma_D\}$.}

\noindent Maka $\Tau$ adalah aljabar-$\sigma$ atas
subhimpunan-subhimpunan $\Bbb R$.
\Prf\ {(i)} $f^{-1}[\emptyset]=\emptyset\in\Sigma_D$, sehingga
$\emptyset\in\Tau$. {(ii)} Jika $E\in\Tau$, maka
$f^{-1}[\Bbb R\setminus E]=D\setminus f^{-1}[E]\in\Sigma_D$,
sehingga $\Bbb R\setminus E\in\Tau$.
{(iii)} Jika $\langle E_n\rangle_{n\in\Bbb N}$ adalah suatu barisan
di dalam $\Tau$, maka
$f^{-1}[\bigcup_{n\in\Bbb N}E_n]
=\bigcup_{n\in\Bbb N}f^{-1}[E_n]\in\Sigma_D$ karena $\Sigma_D$
adalah aljabar-$\sigma$, sehingga $\bigcup_{n\in\Bbb N}E_n\in\Tau$.\ \Qed

Selanjutnya, $\Tau$ memuat semua himpunan berbentuk
$H_a=\ooint{-\infty,a}$ untuk $a\in\Bbb R$, menurut definisi
keterukuran $f$. Hasilnya mengikuti dari argumen yang sudah digunakan
dalam 114G di atas. Pertama, semua subhimpunan terbuka dari $\Bbb R$
termasuk dalam $\Tau$. \Prf\ Misalkan $G\subseteq\Bbb R$ terbuka.
Misalkan $K\subseteq\Bbb Q\times\Bbb Q$ adalah himpunan pasangan
$(q,q')$ bilangan rasional sedemikian sehingga
$\coint{q,q'}\subseteq G$. $K$ terhitung. Selain itu, setiap
$\coint{q,q'}$ termasuk dalam $\Tau$, karena merupakan
$H_{q'}\setminus H_q$. Jadi
$G'=\bigcup_{(q,q')\in K}\coint{q,q'}\in\Tau$.

Menurut definisi $K$, $G'\subseteq G$. Di pihak lain, jika $x\in G$,
terdapat $\delta>0$ sedemikian sehingga
$\ooint{x-\delta,x+\delta}\subseteq G$. Sekarang terdapat bilangan
rasional $q\in\ocint{x-\delta,x}$ dan
$q'\in\ocint{x,x+\delta}$, sehingga $(q,q')\in K$ dan
$x\in\coint{q,q'}\subseteq G'$. Karena $x$ sebarang, $G=G'$ dan
$G\in\Tau$.\ \Qed

Akhirnya, $\Tau$ adalah aljabar-$\sigma$ atas subhimpunan-subhimpunan
$\Bbb R$ yang memuat keluarga himpunan terbuka, sehingga harus memuat
setiap himpunan Borel, menurut definisi himpunan Borel (111G).

\medskip

{\bf (g)} Jika $a\in\Bbb R$, maka $\{y:h(y)<a\}$ berbentuk
$E\cap\dom h$, dengan $E$ suatu subhimpunan Borel dari $\Bbb R$.
Selanjutnya, $f^{-1}[E]$ berbentuk $F\cap\dom f$, dengan $F\in\Sigma$,
menurut (f) di atas. Jadi

\Centerline{$\{x:(hf)(x)<a\}=F\cap\dom hf\in\Sigma_{\dom hf}$.}

\noindent Karena $a$ sebarang, $hf$ terukur.

\medskip

{\bf (h)} Intinya ialah bahwa
$\Sigma_{A\cap\dom f}=\{E\cap A:E\in\Sigma_{\dom f}\}$. Jadi, jika
$a\in\Bbb R$,

\Centerline{$\{x:(f\restr
A)(x)<a\}=A\cap\{x:f(x)<a\}\in\Sigma_{\dom(f\restr A)}$.}
}%end of proof of 121E

\cmmnt{\medskip

\noindent{\bf Catatan} Tentu saja, bagian (c) dari teorema ini hanya
merupakan hasil menggabungkan (a) dan (d), sedangkan (e) merupakan
akibat dari (d), (g), dan fakta bahwa fungsi kontinu terukur Borel
(121Db).

Saya harap Anda mengenali teknik dalam bukti bagian (d) sebagai suatu
versi argumen yang dapat digunakan untuk membuktikan bahwa limit hasil
kali adalah hasil kali limit, atau bahwa hasil kali fungsi-fungsi
kontinu bersifat kontinu. Sebenarnya, (b) dan (d) di sini, bersama
dengan teorema-teorema tentang jumlah dan hasil kali limit, merupakan
akibat dari fakta bahwa penjumlahan dan perkalian adalah fungsi
kontinu. Dalam 121K saya memberikan suatu hasil umum yang dapat
digunakan untuk memanfaatkan fakta-fakta semacam itu.

Sesungguhnya, bagian (f) di sini merupakan inti konsep fungsi bernilai
real yang `terukur'. Inti definisi dalam 121B-121C ialah bahwa
aljabar-$\sigma$ Borel pada $\Bbb R$ dapat dibangkitkan oleh salah satu
di antara keluarga-keluarga
$\{\ooint{-\infty,a}:a\in\Bbb R\}$,
$\{\ocint{-\infty,a}:a\in\Bbb R\},\ldots$. (Lihat 121Yc(ii).)
Ada banyak cara membahas pokok ini dengan jauh lebih ringkas daripada
yang saya lakukan, dengan konsekuensi tingkat abstraksi yang sedikit
lebih tinggi.
}%end of comment


\leader{121F}{Teorema} Misalkan $X$ suatu himpunan dan $\Sigma$ suatu
aljabar-$\sigma$ atas subhimpunan-subhimpunan $X$. Misalkan
$\langle f_n\rangle_{n\in\Bbb N}$ suatu barisan fungsi bernilai real
yang $\Sigma$-terukur, dengan domain-domain yang termuat dalam $X$.

(a) Definisikan suatu fungsi $\lim_{n\to\infty}f_n$ dengan menuliskan

\Centerline{$(\lim_{n\to\infty}f_n)(x)=\lim_{n\to\infty}f_n(x)$}

\noindent untuk semua
$x\in\bigcup_{n\in\Bbb N}\bigcap_{m\ge n}\dom f_m$ yang limitnya ada
dalam $\Bbb R$. Maka $\lim_{n\to\infty}f_n$ merupakan fungsi
$\Sigma$-terukur.


(b) Definisikan suatu fungsi $\sup_{n\in\Bbb N}f_n$ dengan menuliskan

\Centerline{$(\sup_{n\in\Bbb N}f_n)(x)=\sup_{n\in\Bbb N}f_n(x)$}

\noindent untuk semua $x\in\bigcap_{n\in\Bbb N}\dom f_n$ yang
supremumnya ada dalam $\Bbb R$. Maka $\sup_{n\in\Bbb N}f_n$ merupakan
fungsi $\Sigma$-terukur.

(c) Definisikan suatu fungsi $\inf_{n\in\Bbb N}f_n$ dengan menuliskan

\Centerline{$(\inf_{n\in\Bbb N}f_n)(x)=\inf_{n\in\Bbb N}f_n(x)$}

\noindent untuk semua $x\in\bigcap_{n\in\Bbb N}\dom f_n$ yang
infimumnya ada dalam $\Bbb R$. Maka $\inf_{n\in\Bbb N}f_n$ merupakan
fungsi $\Sigma$-terukur.

(d) Definisikan suatu fungsi $\limsup_{n\to\infty}f_n$ dengan
menuliskan

\Centerline{$(\limsup_{n\to\infty}f_n)(x)
=\limsup_{n\to\infty}f_n(x)$}

\noindent untuk semua
$x\in\bigcup_{n\in\Bbb N}\bigcap_{m\ge n}\dom f_m$ yang $\limsup$-nya
ada dalam $\Bbb R$. Maka $\limsup_{n\to\infty}f_n$ merupakan fungsi
$\Sigma$-terukur.

(e) Definisikan suatu fungsi $\liminf_{n\to\infty}f_n$ dengan
menuliskan

\Centerline{$(\liminf_{n\to\infty}f_n)(x)
=\liminf_{n\to\infty}f_n(x)$}

\noindent untuk semua
$x\in\bigcup_{n\in\Bbb N}\bigcap_{m\ge n}\dom f_m$ yang $\liminf$-nya
ada dalam $\Bbb R$. Maka $\liminf_{n\to\infty}f_n$ merupakan fungsi
$\Sigma$-terukur.

\proof{ Untuk $n\in\Bbb N$, $a\in\Bbb R$, pilih
$H_n(a)\in\Sigma$ sedemikian sehingga
$\{x:f_n(x)\le a\}=H_n(a)\cap \dom f_n$. Sekarang, bukti-buktinya
cukup dengan mengamati fakta-fakta berikut:

\medskip

{\bf (a)}
$\{x:(\lim_{n\to\infty}f_n)(x)\le a\}
=\dom(\lim_{n\to\infty}f_n)
  \cap\bigcap_{k\in\Bbb N}\bigcup_{n\in\Bbb N}
    \bigcap_{m\ge n}H_m(a+2^{-k})$;

\medskip


{\bf (b)} $\{x:(\sup_{n\in\Bbb N}f_n)(x)\le a\}
=\dom(\sup_{n\in\Bbb N}f_n)\cap\bigcap_{n\in\Bbb N}H_n(a)$;

\medskip

{\bf (c)} $\inf_{n\in\Bbb N}f_n=-\sup_{n\in\Bbb N}(-f_n)$;

\medskip

{\bf (d)} $\limsup_{n\to\infty}f_n
=\lim_{n\to\infty}\sup_{m\in\Bbb N}f_{m+n}$;

\medskip

{\bf (e)} $\liminf_{n\to\infty}f_n
=-\limsup_{n\to\infty}(-f_n)$.
}%end of proof of 121F


\cmmnt{
\leader{121G}{Catatan} Pada titik inilah kita untuk pertama kalinya
berhadapan secara nyata dengan persoalan fungsi yang tidak terdefinisi
di seluruh ruang. (Hasil bagi $f/g$ dalam 121Ee juga tidak harus
terdefinisi di seluruh domain bersama $f$ dan $g$, tetapi hal ini kurang
penting dan lebih mudah ditangani.) Sejak Lebesgue, inti teori ukuran
dan integrasi ialah bahwa kita dapat menangani limit barisan fungsi;
dan himpunan tempat $\lim_{n\to\infty}f_n(x)$ ada dapat sangat tidak
teratur, bahkan bagi fungsi-fungsi $f_n$ yang tampaknya berperilaku
baik. (Jika Anda pernah menjumpai teori deret Fourier, contoh yang tepat
untuk dipikirkan ialah barisan jumlah parsial
$f_n(x)={1\over 2}a_0+\sum_{k=1}^n(a_k\cos kx+b_k\sin kx)$ dari suatu
deret Fourier dengan $\sum_{k=1}^{\infty}(|a_k|+|b_k|)=\infty$, sehingga
deret tersebut tidak dapat dijumlahkan mutlak secara seragam, tetapi
mungkin dapat dijumlahkan secara bersyarat pada titik-titik tertentu.)

Saya telah berusaha menjelaskan domain mana yang saya maksud bagi
fungsi-fungsi $\sup_{n\in\Bbb N}f_n$,
$\lim_{n\to\infty}f_n$, dan seterusnya. Asas penuntunnya ialah bahwa
domain-domain itu harus menjadi himpunan semua $x\in X$ yang untuknya
rumus pendefinisi $\sup_{n\in\Bbb N}f_n(x)$,
$\lim_{n\to\infty}f_n(x)$ dapat ditafsirkan sebagai bilangan real.
(Seperti saya catat dalam 121C, untuk sementara saya menghindari
`$\infty$' sebagai nilai fungsi, walaupun hal itu hanya menimbulkan
sedikit kesulitan dan beberapa rumus lebih alami ditafsirkan dengan
mengizinkannya.) Namun, dalam kasus $\lim$, $\limsup$, $\liminf$,
perlu dicatat bahwa saya menggunakan definisi yang membatasi, yakni
$\lim_{n\to\infty}a_n$ dapat dianggap ada hanya jika terdapat suatu
$n\in\Bbb N$ sedemikian sehingga $a_m$ terdefinisi untuk setiap
$m\ge n$. Ada kalanya lebih alami untuk menerima limit ketika kita
hanya tahu bahwa $a_m$ terdefinisi untuk tak hingga banyak nilai $m$;
tetapi kesepakatan semacam itu dapat membuat 121Fa salah, kecuali jika
diterapkan dengan hati-hati.

Seperti dalam 111E-111F, kita dapat menggunakan gagasan pada bagian (b),
(c) di sini untuk membahas fungsi berbentuk $\sup_{k\in K}f_k$,
$\inf_{k\in K}f_k$ bagi sembarang keluarga
$\langle f_k\rangle_{k\in K}$ fungsi terukur yang diindeks oleh
himpunan terhitung tak kosong $K$.

Dalam teorema ini dan teorema sebelumnya, fungsi-fungsi $f$, $g$,
$f_n$ diizinkan mempunyai domain sembarang, sehingga tidak ada yang
dapat dikatakan tentang domain fungsi-fungsi yang dikonstruksi. Namun,
tentu saja, jika fungsi-fungsi semula mempunyai domain terukur, maka
demikian pula fungsi-fungsi yang dikonstruksi darinya menurut aturan
yang saya usulkan. Saya menjabarkan rinciannya dalam proposisi berikut.
}%end of comment

\leader{121H}{Proposisi} Misalkan $X$ suatu himpunan dan $\Sigma$
suatu aljabar-$\sigma$ atas subhimpunan-subhimpunan $X$; misalkan
$f$, $g$, dan $f_n$, untuk $n\in\Bbb N$, adalah fungsi-fungsi bernilai
real yang $\Sigma$-terukur dan domainnya termasuk dalam $\Sigma$. Maka
semua fungsi

\Centerline{$f+g$, \quad$f\times g$, \quad$f/g$,}

\Centerline{$\sup_{n\in\Bbb N}f_n$, \quad$\inf_{n\in\Bbb N}f_n$,
\quad$\lim_{n\to\infty}f_n$, \quad$\limsup_{n\to\infty}f_n$,
\quad$\liminf_{n\to\infty}f_n$}

\noindent mempunyai domain yang termasuk dalam $\Sigma$. Selain itu,
jika $h$ adalah fungsi bernilai real terukur Borel yang didefinisikan
pada suatu subhimpunan Borel dari $\Bbb R$, maka $\dom hf\in\Sigma$.

\proof{ Untuk dua fungsi pertama, kita mempunyai
$\dom(f+g)=\dom(f\times g)=\dom f\cap \dom g$. Selanjutnya, jika $E$
adalah suatu subhimpunan Borel dari $\Bbb R$, terdapat $H\in\Sigma$
sedemikian sehingga $f^{-1}[E]=H\cap\dom f$; jadi
$f^{-1}[E]\in\Sigma$. Dengan demikian

\Centerline{$\dom hf=f^{-1}[\dom h]\in\Sigma$.}

\noindent Dengan menetapkan $h(a)=1/a$ untuk
$a\in\Bbb R\setminus\{0\}$, kita melihat bahwa
$\dom(1/f)\in\Sigma$. ($\dom h=\Bbb R\setminus\{0\}$ adalah himpunan
Borel karena terbuka.) Demikian pula, $\dom(1/g)$ dan
$\dom(f/g)=\dom f\cap\dom(1/g)$ termasuk dalam $\Sigma$.

Sekarang tinjau kombinasi-kombinasi tak hingga. Tetapkan
$H_n(a)=\{x:x\in\dom f_n,\,f_n(x)<a\}$ untuk
$n\in\Bbb N$, $a\in\Bbb R$; seperti baru saja dijelaskan, setiap
$H_n(a)$ termasuk dalam $\Sigma$. Sekarang

\Centerline{$\dom(\sup_{n\in\Bbb N}f_n)
=\bigcup_{m\in\Bbb N}\bigcap_{n\in\Bbb N}H_n(m)\in\Sigma$.}

\noindent Selanjutnya, $|f_m-f_n|$ terukur, dengan domain yang termasuk
dalam $\Sigma$, untuk semua $m$, $n\in\Bbb N$ (dengan menerapkan
hasil-hasil di atas pada $-f_n=-1\cdot f_n$,
$f_m-f_n=f_m+(-f_n)$, dan $|f_m-f_n|=|\,\,|\circ(f_m-f_n)$), sehingga

\Centerline{$G_{mnk}
=\{x:x\in\dom f_m\cap\dom f_n,\,|f_m(x)-f_n(x)|\le 2^{-k}\}\in\Sigma$}

\noindent untuk semua $m$, $n$, $k\in\Bbb N$. Dengan demikian

\Centerline{$\dom(\lim_{n\to\infty}f_n)
=\{x:\Exists n,\,\langle f_m(x)\rangle_{m\ge n}$ merupakan barisan Cauchy$\}
=\bigcap_{k\in\Bbb N}\bigcup_{n\in\Bbb N}\bigcap_{m\ge n}G_{mnk}
\in\Sigma$.}

Dengan memanipulasi hasil-hasil di atas seperti dalam (c), (d), dan (e)
dari bukti 121F, kita dengan mudah menyelesaikan bukti.
}%end of proof of 121H

\cmmnt{\medskip

\noindent{\bf Catatan} Perhatikan penggunaan Asas Umum Konvergensi
dalam bukti di atas. Saya tidak yakin apakah hal ini akan terasa
`alami' bagi Anda, dan ada metode-metode alternatif; tetapi rumus

\Centerline{$\{x:\lim_{n\to\infty}f_n(x)$ ada dalam $\Bbb R\}
=\{x:\sequencen{f_n(x)}$ merupakan barisan Cauchy$\}$}

\noindent patut disimpan dalam ingatan jangka panjang Anda.
}%end of comment

\leader{*121I}{}\cmmnt{ Saya mengakhiri bagian ini dengan dua hasil
yang dapat dilewati dengan aman pada pembacaan pertama, tetapi yang pada
suatu saat perlu Anda kuasai jika ingin melangkah lebih jauh dalam teori
ukuran daripada bab ini, karena keduanya merupakan bagian-bagian hakiki
dari konsep `fungsi terukur'.

\medskip

\noindent}{\bf Proposisi} Misalkan $X$ suatu himpunan dan $\Sigma$
suatu aljabar-$\sigma$ atas subhimpunan-subhimpunan $X$. Misalkan $D$
suatu subhimpunan $X$ dan $f:D\to\Bbb R$ suatu fungsi. Maka $f$ terukur
jika dan hanya jika terdapat fungsi terukur $h:X\to\Bbb R$ yang
memperluas $f$.

\proof{{\bf (a)} Jika $h:X\to\Bbb R$ terukur dan $f=h\restr D$, maka
$f$ terukur menurut 121Eh.

\medskip

{\bf (b)} Sekarang andaikan bahwa $f$ terukur.

\medskip

\quad{\bf (i)} Untuk setiap $q\in\Bbb Q$, himpunan
$D_q=\{x:x\in D,\,f(x)\le q\}$ termasuk dalam aljabar-$\sigma$
subruang $\Sigma_D$, yakni, terdapat $E_q\in\Sigma$ sedemikian sehingga
$D_q=E_q\cap D$. Tetapkan

\Centerline{$F=X\setminus\bigcup_{q\in\Bbb Q}E_q$,}

\Centerline{$G=\bigcap_{n\in\Bbb N}\bigcup_{q\in\Bbb Q,q\le -n}E_q$;}

\noindent maka $F$ dan $G$ keduanya termasuk dalam $\Sigma$ dan saling
lepas dari $D$. \Prf\ Jika $x\in D$, terdapat $q\in\Bbb Q$ sedemikian
sehingga $f(x)\le q$, sehingga $x\in E_q$ dan $x\notin F$. Selain itu,
terdapat $n\in\Bbb N$ sedemikian sehingga $f(x)>-n$, sehingga
$x\notin E_{q'}$ untuk $q'\le-n$ dan $x\notin G$.\ \Qed

Tetapkan $H=X\setminus(F\cup G)\in\Sigma$. Untuk $x\in H$,

\Centerline{$\{q:q\in\Bbb Q,\,x\in E_q\}$}

\noindent tak kosong dan berbatas bawah, sehingga kita dapat menetapkan

\Centerline{$h(x)=\inf\{q:x\in E_q\}$;}

\noindent untuk $x\in F\cup G$, tetapkan $h(x)=0$. Ini mendefinisikan
$h:X\to\Bbb R$.

\medskip

\quad{\bf (ii)} $h(x)=f(x)$ untuk $x\in D$. \Prf\ Seperti telah
dicatat di atas, $x\in H$. Jika $q\in\Bbb Q$ dan $x\in E_q$, maka
$f(x)\le q$; akibatnya $h(x)\ge f(x)$. Di pihak lain, untuk setiap
$\epsilon>0$, terdapat $q\in\Bbb Q\cap[f(x),f(x)+\epsilon]$, dan
sekarang $x\in E_q$, sehingga $h(x)\le q\le f(x)+\epsilon$; karena
$\epsilon$ sebarang, $h(x)\le f(x)$.\ \Qed

\medskip

\quad{\bf (iii)} $h$ terukur. \Prf\ Jika $a>0$, maka

\Centerline{$\{x:h(x)<a\}=(H\cap\bigcup_{q<a}E_q)\cup(F\cup
G)\in\Sigma$,}

\noindent sedangkan jika $a\le 0$,

\Centerline{$\{x:h(x)<a\}=H\cap\bigcup_{q<a}E_q\in\Sigma$.\ \Qed}

\noindent Ini menyelesaikan bukti.
}%end of proof of 121I

\leader{*121J}{}\cmmnt{ Proposisi berikut dapat memperjelas 121E,
sekaligus sangat diperlukan bagi pekerjaan dalam Jilid 2. Saya mulai
dengan suatu deskripsi berguna tentang himpunan-himpunan Borel dalam
$\BbbR^r$.

\medskip

\noindent}{\bf Lema} Misalkan $r\ge 1$ suatu bilangan bulat, dan
tuliskan $\Cal J$ untuk keluarga subhimpunan-subhimpunan $\BbbR^r$
yang berbentuk $\{x:\xi_i\le\alpha\}$, dengan $i\le r$,
$\alpha\in\Bbb R$, serta $x=(\xi_1,\ldots,\xi_r)$, seperti dalam
\S115. Maka aljabar-$\sigma$ atas subhimpunan-subhimpunan $\BbbR^r$
yang dibangkitkan oleh $\Cal J$ tepat sama dengan aljabar-$\sigma$
$\Cal B$ dari subhimpunan-subhimpunan Borel $\BbbR^r$.

\proof{{\bf (a)} Semua himpunan dalam $\Cal J$ tertutup, sehingga
harus termasuk dalam $\Cal B$; dengan menuliskan $\Sigma$ untuk
aljabar-$\sigma$ yang dibangkitkan oleh $\Cal J$, kita harus mempunyai
$\Sigma\subseteq\Cal B$.

\medskip


{\bf (b)} Langkah berikutnya ialah mengamati bahwa semua interval
setengah terbuka berbentuk

\Centerline{$\ocint{a,b}
=\{x:\alpha_i<\xi_i\le\beta_i\Forall i\le r\}$}

\noindent termasuk dalam $\Sigma$; hal ini karena

\Centerline{$\ocint{a,b}=\bigcap_{i\le
r}(\{x:\xi_i\le\beta_i\}\setminus\{x:\xi_i\le\alpha_i\})$.}

\noindent Akibatnya, semua himpunan terbuka termasuk dalam $\Sigma$.
\Prf\ (Bandingkan bukti 121Ef.) Misalkan $G\subseteq\BbbR^r$ suatu
himpunan terbuka. Misalkan $\Cal I$ adalah himpunan semua interval
berbentuk $\ocint{q,q'}$ yang termuat dalam $G$, dengan
$q$, $q'\in\BbbQ^r$. Maka $\Cal I$ adalah subhimpunan terhitung dari
$\Sigma$, sehingga (karena $\Sigma$ adalah aljabar-$\sigma$)
$\bigcup\Cal I\in\Sigma$. Menurut definisi $\Cal I$,
$\bigcup\Cal I\subseteq G$. Namun, jika $x\in G$, terdapat
$\delta>0$ sedemikian sehingga bola terbuka $U(x,\delta)$ yang berpusat
di $x$ dan berjari-jari $\delta$ termuat dalam $G$ (1A2A).
Sekarang, untuk setiap $i\le r$, kita dapat menemukan bilangan-bilangan
rasional $\alpha_i$, $\beta_i$ sedemikian sehingga

\Centerline{$\xi_i-\Bover{\delta}r
\le\alpha_i<\xi_i\le\beta_i<\xi_i+\Bover{\delta}r$,}

\noindent sehingga

\Centerline{$x\in\ocint{a,b}\subseteq U(x,\delta)\subseteq G$}

\noindent dan $x\in\ocint{a,b}\in\Cal I$. Jadi
$x\in\bigcup\Cal I$. Karena $x$ sebarang,
$G\subseteq\bigcup\Cal I$ dan $G=\bigcup\Cal I\in\Sigma$.\ \Qed

\medskip

{\bf (c)} Dengan demikian, $\Sigma$ adalah aljabar-$\sigma$ himpunan
yang memuat setiap himpunan terbuka, dan harus memuat $\Cal B$,
aljabar-$\sigma$ terkecil semacam itu.
}%end of proof of 121J

\cmmnt{\medskip

\noindent{\bf Catatan} Bandingkan bukti 115G.
}%end of comment

\leader{*121K}{Proposisi} Misalkan $X$ suatu himpunan dan $\Sigma$
suatu aljabar-$\sigma$ atas subhimpunan-subhimpunan $X$. Misalkan
$r\ge 1$ suatu bilangan bulat, dan $f_1,\ldots,f_r$ fungsi-fungsi
terukur yang didefinisikan pada subhimpunan-subhimpunan $X$. Tetapkan
$D=\bigcap_{i\le r}\dom f_i$, dan untuk $x\in D$ tetapkan
$f(x)=(f_1(x),\ldots,f_r(x))\in\BbbR^r$. Maka

(a) untuk setiap himpunan Borel $E\subseteq\BbbR^r$, $f^{-1}[E]$
termasuk dalam aljabar-$\sigma$ subruang $\Sigma_D$;

(b) jika $h$ adalah suatu fungsi terukur Borel dari suatu subhimpunan
$\dom h$ dari $\BbbR^r$ ke $\Bbb R$, maka komposisi $hf$ terukur.

\proof{{\bf (a)(i)} Tinjau himpunan

\Centerline{$\Tau=\{E:E\subseteq\BbbR^r,\,f^{-1}[E]\in\Sigma_D\}$.}

\noindent Maka $\Tau$ adalah aljabar-$\sigma$ atas
subhimpunan-subhimpunan $\BbbR^r$.
\Prf\ (Bandingkan 121Ef.) {\bf ($\pmb{\alpha}$)}
$f^{-1}[\emptyset]=\emptyset\in\Sigma_D$, sehingga
$\emptyset\in\Tau$. {\bf ($\pmb{\beta}$)} Jika $E\in\Tau$, maka
$f^{-1}[\BbbR^r\setminus E]=D\setminus f^{-1}[E]\in\Sigma_D$,
sehingga $\BbbR^r\setminus E\in\Tau$.
{\bf ($\pmb{\gamma}$)} Jika $\langle E_n\rangle_{n\in\Bbb N}$ adalah
suatu barisan di dalam $\Tau$, maka
$f^{-1}[\bigcup_{n\in\Bbb N}E_n]
=\bigcup_{n\in\Bbb N}f^{-1}[E_n]\in\Sigma_D$ karena $\Sigma_D$
adalah aljabar-$\sigma$, sehingga $\bigcup_{n\in\Bbb N}E_n\in\Tau$.\ \Qed

\medskip

\quad{\bf (ii)} Selanjutnya, untuk setiap $i\le r$ dan
$\alpha\in\Bbb R$, $J=\{x:\xi_i\le\alpha\}$ termasuk dalam $\Tau$,
karena

\Centerline{$f^{-1}[J]
=\{x:x\in D,\, f_i(x)\le \alpha\}\in\Sigma_D$.}

\noindent Jadi $\Tau$ memuat keluarga $\Cal J$ dari 121J, dan karena
itu memuat aljabar-$\sigma$ $\Cal B$ yang dibangkitkan oleh $\Cal J$,
yakni, memuat setiap subhimpunan Borel dari $\BbbR^r$.

\medskip

{\bf (b)} Sekarang sisanya mengikuti dari argumen 121Eg. Jika
$a\in\Bbb R$, maka $\{y:y\in\dom h,\,h(y)<a\}$ berbentuk
$E\cap\dom h$, dengan $E$ suatu subhimpunan Borel dari $\BbbR^r$,
sehingga $\{x:x\in\dom(hf),\,(hf)(x)<a\}
=f^{-1}[E]\cap\dom(hf)$ termasuk dalam $\Sigma_{\dom hf}$.
}%end of proof of 121K

\exercises{
\leader{121X}{Latihan dasar $\pmb{>}$(a)}
%\spheader 121Xa
Misalkan $X$ suatu himpunan, $\Sigma$ suatu aljabar-$\sigma$ atas
subhimpunan-subhimpunan $X$, dan $D\subseteq X$. Misalkan
$\sequencen{D_n}$ suatu partisi $D$ menjadi himpunan-himpunan yang
terukur relatif, dan $\sequencen{f_n}$ suatu barisan fungsi bernilai
real terukur sedemikian sehingga $D_n\subseteq\dom f_n$ untuk setiap
$n$. Definisikan $f:D\to\Bbb R$ dengan menetapkan $f(x)=f_n(x)$ setiap
kali $n\in\Bbb N$, $x\in D_n$. Tunjukkan bahwa $f$ terukur.
%121C

\spheader 121Xb Misalkan $X$ suatu himpunan dan $\Sigma$ suatu
aljabar-$\sigma$ atas subhimpunan-subhimpunan $X$. Jika $f$ dan $g$
adalah fungsi-fungsi bernilai real terukur yang didefinisikan pada
subhimpunan-subhimpunan $X$, tunjukkan bahwa $f^+$, $f^-$,
$f\wedge g$, dan $f\vee g$ terukur, dengan

\Centerline{$f^+(x)=\max(f(x),0)$ untuk $x\in\dom f$,}

\Centerline{$f^-(x)=\max(-f(x),0)$ untuk $x\in\dom f$,}

\Centerline{$(f\vee g)(x)=\max(f(x),g(x))$ untuk
$x\in\dom f\cap\dom g$,}

\Centerline{$(f\wedge g)(x)=\min(f(x),g(x))$ untuk
$x\in\dom f\cap\dom g$.}
%121E

\sqheader 121Xc Misalkan $(X,\Sigma,\mu)$ suatu ruang ukur. Tuliskan
$\eusm L^0$ untuk himpunan fungsi bernilai real $f$ sedemikian sehingga
($\alpha$) $\dom f$ adalah subhimpunan koterabaikan dari $X$
($\beta$) terdapat himpunan koterabaikan $E\subseteq X$ sedemikian
sehingga $f\restr E$ terukur.
(i) Tunjukkan bahwa himpunan $E$ pada klausa ($\beta$) dalam kalimat
terakhir dapat diambil termasuk dalam $\Sigma$ dan termuat dalam
$\dom f$. (ii) Tunjukkan bahwa jika $f$, $g\in\eusm L^0$ dan
$c\in\Bbb R$, maka $f+g$, $cf$, $f\times g$, $|f|$, $f^+$, $f^-$,
$f\wedge g$, $f\vee g$ semuanya termasuk dalam $\eusm L^0$.
(iii) Tunjukkan bahwa jika $f$, $g\in\eusm L^0$ dan $g\ne 0$ a.e.,
maka $f/g\in\eusm L^0$. (iv) Tunjukkan bahwa jika
$\sequencen{f_n}$ adalah suatu barisan di dalam $\eusm L^0$, maka
fungsi-fungsi

\Centerline{$\lim_{n\to\infty}f_n$,\quad$\sup_{n\in\Bbb N}f_n$,
\quad$\inf_{n\in\Bbb N}f_n$,\quad $\limsup_{n\to\infty}f_n$,
\quad$\liminf_{n\to\infty}f_n$}

\noindent termasuk dalam $\eusm L^0$ setiap kali fungsi-fungsi itu
terdefinisi hampir di mana-mana sebagai fungsi bernilai real.
(v) Tunjukkan bahwa jika $f\in\eusm L^0$ dan
$h:\Bbb R\to\Bbb R$ terukur Borel, maka $hf\in\eusm L^0$.
%121F

\sqheader 121Xd Tinjau empat keluarga subhimpunan $\Bbb R$ berikut:

\Centerline{$\Cal A_1=\{\ooint{-\infty,a}:a\in\Bbb R\}$,
\quad$\Cal A_2=\{\ocint{-\infty,a}:a\in\Bbb R\}$,}

\Centerline{$\Cal A_3=\{\ooint{a,\infty}:a\in\Bbb R\}$,
\quad$\Cal A_4=\{\coint{a,\infty}:a\in\Bbb R\}$.}

\noindent Tunjukkan bahwa untuk setiap $j$, aljabar-$\sigma$ atas
subhimpunan-subhimpunan $\Bbb R$ yang dibangkitkan oleh $\Cal A_j$
adalah aljabar-$\sigma$ himpunan-himpunan Borel.
%121J

\spheader 121Xe Misalkan $D$ sembarang subhimpunan $\BbbR^r$, dengan
$r\ge 1$. Tuliskan $\frak T_D$ untuk himpunan
$\{G\cap D:G\subseteq\BbbR^r$ merupakan himpunan terbuka$\}$.
(i) Tunjukkan bahwa $\frak T_D$ memenuhi sifat-sifat himpunan terbuka
yang tercantum dalam 1A2B. (ii) Misalkan $\Cal B$ aljabar-$\sigma$
himpunan-himpunan Borel dalam $\BbbR^r$, dan $\Cal B(D)$
aljabar-$\sigma$ subruang pada $D$. Tunjukkan bahwa $\Cal B(D)$ tepat
sama dengan aljabar-$\sigma$ atas subhimpunan-subhimpunan $D$ yang
dibangkitkan oleh $\frak T_D$. \Hint{($\alpha$) amati bahwa
$\frak T_D\subseteq\Cal B(D)$ ($\beta$) tinjau
$\{E:E\subseteq\BbbR^r,\,E\cap D$ termasuk dalam aljabar-$\sigma$
yang dibangkitkan oleh $\frak T_D\}$.}
%121K

\spheader 121Xf Misalkan $(X,\Sigma,\mu)$ suatu ruang ukur dan
definisikan $\eusm L^0$ seperti dalam 121Xc. Tunjukkan bahwa jika
$f_1,\ldots,f_r$ termasuk dalam $\eusm L^0$ dan
$h:\BbbR^r\to\Bbb R$ terukur Borel, maka $h(f_1,\ldots,f_r)$
termasuk dalam $\eusm L^0$.
%121Xc 121K

\leader{121Y}{Latihan lanjutan (a)}
%\spheader 121Ya
Misalkan $X$ dan $Y$ himpunan-himpunan, $\Sigma$ suatu
aljabar-$\sigma$ atas subhimpunan-subhimpunan $X$, $\phi:X\to Y$ suatu
fungsi, dan $g$ suatu fungsi bernilai real yang didefinisikan pada suatu
subhimpunan $Y$. Tetapkan
$\Tau=\{F:F\subseteq Y,\,\phi^{-1}[F]\in\Sigma\}$; maka $\Tau$ adalah
aljabar-$\sigma$ atas subhimpunan-subhimpunan $Y$ (lihat 111Xc).
(i) Tunjukkan bahwa jika $g$ $\Tau$-terukur, maka $g\phi$
$\Sigma$-terukur.
(ii) Berikan suatu contoh ketika $g\phi$ $\Sigma$-terukur, tetapi $g$
tidak $\Tau$-terukur.
(iii) Tunjukkan bahwa jika $g\phi$ $\Sigma$-terukur dan
{\it salah satu} dari syarat berikut berlaku: $\phi$ injektif,
{\it atau} $\dom(g\phi)\in\Sigma$, {\it atau}
$\phi[X]\subseteq\dom g$, maka $g$ $\Tau$-terukur.\footnote{Saya
berterima kasih kepada P. Wallace Thompson karena telah menunjukkan
kekeliruan dalam versi asli latihan ini.}
%121C

\spheader 121Yb Misalkan $X$ dan $Y$ himpunan-himpunan, $\Tau$ suatu
aljabar-$\sigma$ atas subhimpunan-subhimpunan $Y$, dan $\phi:X\to Y$
suatu fungsi. Tetapkan $\Sigma=\{\phi^{-1}[F]:F\in\Tau\}$, seperti
dalam 111Xd. Tunjukkan bahwa suatu fungsi $f:X\to\Bbb R$
$\Sigma$-terukur jika dan hanya jika terdapat fungsi $\Tau$-terukur
$g:Y\to\Bbb R$ sedemikian sehingga $f=g\phi$.
%121I

\spheader 121Yc Misalkan $X$ dan $Y$ himpunan-himpunan, dan $\Sigma$,
$\Tau$ masing-masing aljabar-$\sigma$ atas subhimpunan-subhimpunan
$X$, $Y$. Suatu fungsi $\phi:X\to Y$ disebut
{\bf $(\Sigma,\Tau)$-terukur} jika $\phi^{-1}[F]\in\Sigma$ untuk
setiap $F\in\Tau$.
(i) Tunjukkan bahwa jika $\Sigma$, $\Tau$, $\Upsilon$ masing-masing
adalah aljabar-$\sigma$ atas subhimpunan-subhimpunan $X$, $Y$, $Z$,
dan $\phi:X\to Y$ $(\Sigma,\Tau)$-terukur,
$\psi:Y\to Z$ $(\Tau,\Upsilon)$-terukur, maka $\psi\phi:X\to Z$
$(\Sigma,\Upsilon)$-terukur.
(ii) Andaikan $\Cal A\subseteq\Tau$ sedemikian sehingga $\Tau$
merupakan aljabar-$\sigma$ atas subhimpunan-subhimpunan $Y$ yang
dibangkitkan oleh $\Cal A$ (111Gb). Tunjukkan bahwa $\phi:X\to Y$
$(\Sigma,\Tau)$-terukur jika dan hanya jika
$\phi^{-1}[A]\in\Sigma$ untuk setiap $A\in\Cal A$.
(iii) Untuk $r\ge 1$, tuliskan $\Cal B_r$ untuk aljabar-$\sigma$
subhimpunan-subhimpunan Borel dari $\BbbR^r$. Tunjukkan bahwa jika $X$
sembarang himpunan dan $\Sigma$ suatu aljabar-$\sigma$ atas
subhimpunan-subhimpunan $X$, maka suatu fungsi $f:X\to\BbbR^r$
$(\Sigma,\Cal B_r)$-terukur jika dan hanya jika
$\pi_if:X\to\Bbb R$ $(\Sigma,\Cal B_1)$-terukur untuk setiap
$i\le r$, dengan $\pi_i(x)=\xi_i$ untuk $i\le r$,
$x=(\xi_1,\ldots,\xi_r)\in\BbbR^r$.
(iv) Rumuskan kembali gagasan-gagasan ini untuk fungsi-fungsi yang
terdefinisi sebagian.
%121+

\spheader 121Yd Misalkan $X$ suatu himpunan dan $\Sigma$ suatu
aljabar-$\sigma$ atas subhimpunan-subhimpunan $X$. Untuk $r\ge 1$,
$D\subseteq X$, katakan bahwa suatu fungsi $\phi:D\to\BbbR^r$
{\bf terukur} jika $\phi^{-1}[G]$ terukur relatif dalam $D$ untuk
setiap himpunan terbuka $G\subseteq\BbbR^r$. Jika $X=\BbbR^s$ dan
$\Sigma$ adalah aljabar-$\sigma$ $\Cal B_s$ dari subhimpunan-subhimpunan
Borel $\BbbR^s$, katakan bahwa $\phi$ {\bf terukur Borel}.
(i) Tunjukkan bahwa $\phi$ terukur dalam pengertian ini jika dan hanya
jika semua fungsi koordinatnya $\phi_i:D\to\Bbb R$ terukur dalam
pengertian 121C, dengan $\phi(x)=(\phi_1(x),\ldots,\phi_r(x))$ untuk
$x\in D$. (Khususnya, definisi ini sesuai dengan 121C ketika $r=1$.)
(ii) Tunjukkan bahwa $\phi:D\to\BbbR^r$ terukur jika dan hanya jika
$(\Sigma,\Cal B_r)$-terukur dalam pengertian 121Yc.
(iii) Tunjukkan bahwa jika $\phi:D\to\BbbR^r$ terukur dan
$\psi:E\to\BbbR^s$ terukur Borel, dengan $E\subseteq\BbbR^r$, maka
$\psi\phi:\phi^{-1}[E]\to\BbbR^s$ terukur.
(iv) Tunjukkan bahwa setiap fungsi kontinu dari suatu subhimpunan
$\BbbR^s$ ke $\BbbR^r$ terukur Borel.
%121Yc

\spheader 121Ye Misalkan $X$ suatu himpunan dan $\theta$ suatu ukuran
luar pada $X$; misalkan $\mu$ ukuran yang didefinisikan dari $\theta$
dengan metode \Caratheodory, dan $\Sigma$ domainnya. Andaikan bahwa
$f:X\to\Bbb R$ adalah suatu fungsi sedemikian sehingga

\Centerline{$\theta\{x:x\in A,\,f(x)\le a\}
  +\theta\{x:x\in A,\,f(x)\ge b\}\le\theta A$}

\noindent setiap kali $A\subseteq X$ dan $a<b$ dalam $\Bbb R$.
Tunjukkan bahwa $f$ $\Sigma$-terukur. ({\it Petunjuk\/}: andaikan
$a\in\Bbb R$ dan $\theta A<\infty$. Tetapkan

\Centerline{$B_k=\{x:x\in A,\,
 a+\Bover1{2k+2}\le f(x)\le a+\Bover1{2k+1}\}$,}

\Centerline{$B'_k=\{x:x\in A,\,
  a+\Bover1{2k+3}\le f(x)\le a+\Bover1{2k+2}\}$}

\noindent untuk $k\in\Bbb N$. Tunjukkan bahwa
$\sum_{k=0}^{\infty}\theta B_k\le\theta A$, dan periksa hasil serupa
untuk $B'_k$. Dari sini, tunjukkan bahwa

\Centerline{$\theta\{x:x\in A,\,f(x)>a\}
=\lim_{k\to\infty}\theta\{x:x\in A,\,f(x)\ge a+\Bover1k\}$.)}
%121+
}%end of exercises

\endnotes{
\Notesheader{121} Dalam bagian ini saya telah memberikan tidak kurang
dari tiga definisi `fungsi terukur', dalam 121C, 121Yc, dan 121Yd.
Sebenarnya, definisi terakhir barangkali yang terpenting dan menjadi
panduan terbaik menuju gagasan-gagasan selanjutnya. Meskipun demikian,
sebagian
besar penerapan merujuk pada fungsi bernilai real, dan empat syarat
ekuivalen dalam 121B adalah yang paling alami dan paling mudah
digunakan. Fakta bahwa semuanya berimpit dengan syarat dalam 121Yd
bersesuaian dengan fakta bahwa semuanya berbentuk

\Centerline{$f^{-1}[E]\in\Sigma_D$ untuk setiap $E\in\Cal A$}

\noindent dengan $\Cal A$ suatu keluarga subhimpunan $\Bbb R$ yang
membangkitkan aljabar-$\sigma$ Borel (121Xd, 121Yc(ii)).

Kelas fungsi terukur mungkin merupakan kelas terluas yang sejauh ini
pernah Anda lihat, jika kita tidak menghitung keluarga semua fungsi
bernilai real; semua fungsi yang mudah dideskripsikan antara
subhimpunan-subhimpunan $\Bbb R$ bersifat terukur. Ruang fungsi
terukur bukan hanya tertutup terhadap penjumlahan, perkalian, dan
komposisi dengan fungsi kontinu (121E), tetapi setiap operasi alami
yang bekerja pada suatu barisan fungsi terukur akan menghasilkan fungsi
terukur (121F, 121Xb, 121Xa). Namun, {\it tidak} selalu benar bahwa
komposisi dua fungsi terukur Lebesgue dari $\Bbb R$ ke dirinya sendiri
terukur Lebesgue; saya memberikan contoh tandingan dalam 134Ib.
Intinya di sini ialah bahwa suatu fungsi disebut `terukur' jika fungsi
tersebut $(\Sigma,\Cal B)$-terukur, dalam bahasa 121Yc, dengan
$\Cal B$ aljabar-$\sigma$ himpunan-himpunan Borel. Fungsi semacam itu
dapat saja gagal menjadi $(\Sigma,\Sigma)$-terukur jika $\Sigma$
memuat $\Cal B$ secara sejati, sehingga argumen alami untuk komposisi
(121Yc(i)) gagal. Meskipun demikian, karena alasan-alasan yang akan
saya singgung dalam \S134, fungsi yang tidak terukur Lebesgue sulit
ditemukan, dan hanya muncul secara alami dalam jenis-jenis analisis real
yang sangat terspesialisasi. Karena itu, Anda dapat mendekati pertanyaan
apakah suatu fungsi tertentu terukur Lebesgue dengan keyakinan yang
wajar bahwa fungsi tersebut memang terukur, dan bahwa buktinya sekadar
menjadi tantangan bagi teknik Anda.

Anda akan mengamati bahwa hasil-hasil 121E sebagian besar tercakup oleh
121I-121K, yang juga mencakup sebagian besar 114G dan 115G; dan bahwa
121Kb tercakup oleh 121Yd(iii). Anda dapat menganggap diri telah
menguasai bagian pokok bahasan ini ketika pemaparan saya terasa
membosankan karena berulang-ulang. Tentu saja, untuk menurunkan 121Ed
dari 121K, misalnya, Anda harus mengetahui bahwa perkalian, dipandang
sebagai fungsi dari $\BbbR^2$ ke $\Bbb R$, bersifat kontinu, sehingga
terukur Borel; buktinya tertanam dalam bukti yang saya berikan untuk
121Ed (lihat rumus untuk $\eta$ di pertengahan bukti).
}%end of notes

\discrpage
